\documentclass[a4paper]{article}
\usepackage[a4paper, margin=2cm]{geometry}
\usepackage[utf8]{inputenc}
\usepackage[T1]{fontenc}
\usepackage[indonesian]{babel}

\begin{document}

\section*{Pengenalan Pengembangan Game dan Unity}
konsep dasar game, genre, pengenalan antarmuka dan proyek Unity

\section*{Dasar Pemrograman C\# di Unity}
struktur skrip, fungsi \texttt{Start()} dan \texttt{Update()}, variabel

\section*{GameObject dan Komponen Unity}
Transform, Rigidbody, Collider, Prefab

\section*{Manajemen Aset dan Scene}
impor aset (grafik, model), pengaturan scene, pencahayaan dasar

\section*{Minggu 5: Kontrol Pemain dan Input}
pemrosesan input (keyboard, mouse), menggerakkan objek

\section*{Minggu 6: Sistem Fisika}
penggunaan Rigidbody, Collider untuk simulasi fisika pada objek (misalnya gravitasi, tumbukan)

\section*{Minggu 7: Animasi Game}
Animator, Animator Controller, dan animasi sprite/3D (transisi animasi)

\section*{Minggu 8: Antarmuka Pengguna (UI)}
Canvas, elemen UI (button, text), membuat HUD

\section*{Minggu 9: Audio dalam Unity}
penggunaan AudioSource, pengaturan musik dan efek suara pada game

\section*{Minggu 10: Kecerdasan Buatan (AI) dasar}
membuat agen sederhana (misalnya musuh) menggunakan NavMesh atau skrip dasar

\section*{Minggu 11: Mekanika Permainan}
sistem skor, kondisi menang/kalah, state machine


\end{document}

